
\pagebreak
\chapter{ASP.NET and Blazor.}


\medskip

ASP.NET is a robust framework developed by Microsoft for building modern web applications. It provides a scalable and high-performance environment for web development, supporting various technologies such as MVC, Razor Pages, Web APIs, and Blazor.

Blazor is a relatively new framework within the ASP.NET ecosystem that enables developers to build interactive web applications using C\# and .NET instead of JavaScript. With Blazor, developers can create both client-side and server-side web applications with a consistent programming model.

Blazor is a component-based UI framework that allows developers to build web applications using C\# and Razor templates instead of JavaScript. It offers two hosting models: Blazor Server (the application runs on the server, and UI updates are sent over a SignalR connection) and Blazor WebAssembly (the application runs entirely in the user's browser using WebAssembly)

ASP.NET and Blazor together provide a powerful framework for building modern web applications. Whether using Blazor Server or WebAssembly, developers can leverage C\# and .NET to create interactive and scalable web applications.


\textbf{Implementation.}

Let's create a new solution using Visual Studio, and choose Blazor Web App from the list of available template options.

After we have created a new solution with Blazor project (in our case, we named the project 'demo-sln'),

When we run our application, we are greeted with a standard landing page, that we can now configure to our liking.


Let's start by adding folder 'Models' and adding new item 'Product.cs' that would encapsulate information about certain product our application would be handling.

\begin{lstlisting}
namespace demo_sln.Models;

public class Product
{
    public int Id { get; set; }
    public string? Name { get; set; }
    public string? Description { get; set; }
    public decimal Price { get; set; }
    public string? Category { get; set; }
}

\end{lstlisting}

Now we have to create a service, that would be handling these products. This service will be returning products upon request and adding new products to a mock database (we will not be actually creating a database as a part of this particular lesson). Let's name our service 'ProductService.cs'. We will first create an interface and then an actual implementation class:

\begin{lstlisting}
using demo_sln.Models;

namespace demo_sln.Services;

public interface IProductsService
{
    void AddProduct(Product product);
    List<Product> GetProducts();
    Product GetProduct(int id);
}

\end{lstlisting}

Implementation of that interface will inherit the interface:

\begin{lstlisting}
using demo_sln.Models;

namespace demo_sln.Services;

public class ProductsService : IProductsService
{
    private List<Product> _products;

    public ProductsService()
    {
        _products = new List<Product>()
            {
                new Product { Id = 1, Name = "Test name", Description = "Descr"},
                new Product { Id = 2, Name = "Test name 2", Description = "Descr 2"}
            };
    }

    public List<Product> GetProducts()
    {
        return _products;
    }

    public void AddProduct(Product product)
    {
        product.Id = _products.Count() + 1;
        _products.Add(product);
    }

    public Product GetProduct(int id)
    {
        var product = _products.FirstOrDefault(p => p.Id == id);
        if (product == null)
        {
            throw new KeyNotFoundException();
        }
        return product;
    }
}

\end{lstlisting}

A few notes about interfaces and why we need them in C\#:

An interface in C\# is a contract that defines a set of methods and properties without providing an implementation. Classes that implement the interface must provide concrete implementations for all its members.

Using interfaces in C\# provides several benefits. First of all, they provide abstraction - they hide implementation details and expose only necessary functionalities. Loose coupling - interfaces reduce dependencies between classes, making the system more modular. When it comes to testability - they make unit testing easier by allowing mock implementations. And code maintainability - simplify updating or replacing components without breaking the existing code.

Here we use interfaces to provide dependency injection to our code: design pattern that allows injecting dependencies instead of creating them inside a class. This approach improves testability and maintainability. Thus, we use at all times interfaces instead of actual implementations.


Here is an example of how we can provide actual dependency injection of our ProductService:

\begin{lstlisting}
builder.Services.AddSingleton<IProductsService, ProductsService>();
\end{lstlisting}

Note, that we choose AddSingleton option - that is important. We will provide more information on singleton pattern later - for now what is important, is that this way we can ensure data consistency for the application lifetime without resorting to actually adding database to our application for now.

Now let's add our first razor component - file '' to folder Components//Pages

\begin{lstlisting}
@using demo_sln.Models

<div class="item-row">
    <p class="item-name">@Product?.Name</p>

    <p class="item-description">@Product?.Description</p>
</div>

@code {
    [Parameter]
    public Product? Product { get; set; }
}

\end{lstlisting}

This is a Razor component. Blazor Razor components are the building blocks of Blazor applications. Think of them as reusable UI pieces, like buttons, forms, or entire pages, written in a mix of C\# and HTML. Here in the 'code' segment we see C\# code inside the component, and '@Product' variable binds data in code part with the view. Product is market with attribute '[Parameter]' which means that it is a parameter that is passed to that component from another component. Before we define the higher level component, let's first add a ProductItem.razor.css file, that would contain styles for this component:

\begin{lstlisting}
.item-row {
    display: flex;
    align-items: center;
    justify-content: space-between;
    padding: 10px;
    border: 1px solid #ddd;
    border-radius: 8px;
    background-color: #f9f9f9;
    box-shadow: 2px 2px 5px rgba(0, 0, 0, 0.1);
    margin: 10px auto;
}

.item-name {
    font-weight: bold;
    font-size: 18px;
    color: #333;
    margin-right: 10px;
}

.item-description {
    font-size: 14px;
    color: #666;
    flex-grow: 1;
}
\end{lstlisting}

And now let's define top-level component Products.razor:

\begin{lstlisting}
@page "/products"
@using demo_sln.Models
@using demo_sln.Services
@rendermode InteractiveServer
@inject IProductsService _productsService

<h3>Products</h3>

@if (_products == null)
{
    <h3>Please wait a bit</h3>
}
else
{
    <div>
        @foreach (var product in _products)
        {
            <ProductItem Product="product"></ProductItem>
        }
    </div>
}
<div>
    <input @bind="_name" />
    <input @bind="_description" />
    <button @onclick="SubmitData">Click me to add new product</button>

</div>



@code {
    private List<Product>? _products;

    private string? _name;
    private string? _description;

    protected override Task OnInitializedAsync()
    {
        _products = _productsService.GetProducts();
        return base.OnInitializedAsync();
    }

    private void SubmitData()
    {
        _productsService.AddProduct(
            new Product() { Name = _name, Description = _description }
        );
        _name = String.Empty;
        _description = String.Empty;
        _products = _productsService.GetProducts();
    }
}

\end{lstlisting}

In this component:
\begin{itemize}

\item @page "/products" - Defines the route where this component will be accessible.

\item @inject IProductsService \_productsService - Injects a service to fetch product data.

\item @if (\_products == null) - Displays a loading message if data isn't available yet.

\item @foreach (var product in \_products) - Iterates over the products list and displays them using the ProductItem component.

\item @bind - Binds input fields to \_name and \_description properties.

\item @onclick="SubmitData" - Calls SubmitData() when the button is clicked.
\end{itemize}


Before we start our application, we need to add a link that would re-address us to the products list - we do it in the NavMenu.razor component in folder Layout, modify it in the following way:



\begin{lstlisting}
<div class="top-row ps-3 navbar navbar-dark">
    <div class="container-fluid">
        <a class="navbar-brand" href="">demo-sln</a>
    </div>
</div>

<input type="checkbox" title="Navigation menu" class="navbar-toggler" />

<div class="nav-scrollable" onclick="document.querySelector('.navbar-toggler').click()">
    <nav class="flex-column">
        <div class="nav-item px-3">
            <NavLink class="nav-link" href="" Match="NavLinkMatch.All">
                <span class="bi bi-house-door-fill-nav-menu" aria-hidden="true"></span> Home
            </NavLink>
        </div>

        <div class="nav-item px-3">
            <NavLink class="nav-link" href="products">
                <span class="bi bi-plus-square-fill-nav-menu" aria-hidden="true"></span> Products
            </NavLink>
        </div>
    </nav>
</div>
\end{lstlisting}

Now suppose we want to add a page for each individual product item. We can do that with routing: define a route pattern with id of the product in the address of the page - add item ProductPage.razor:

\begin{lstlisting}
@page "/products/{Id:int}"
@inject IProductsService _productsService
@using demo_sln.Models
@using demo_sln.Services
<h3>Product</h3>

<p>@_product?.Name</p>

<p>@_product?.Description</p>


@code {
    [Parameter]
    public int Id { get; set; }

    private Product? _product { get; set; }

    protected override Task OnInitializedAsync()
    {
        _product = _productsService.GetProduct(Id);
        return base.OnInitializedAsync();
    }
}

\end{lstlisting}

Here directive @page "/products/{Id:int}" defines the route by which we can access that component.
Now we modify our ProductItem.razor component a bit to have a link to that page:

\begin{lstlisting}
@using demo_sln.Models

<div class="item-row">
    <span class="item-name">@Product?.Name</span>

    <span class="item-description">@Product?.Description</span>

    <NavLink href=@link>Click</NavLink>
</div>

@code {
    [Parameter]
    public Product? Product { get; set; }

    public string link => "/products/" + Product?.Id;
}
\end{lstlisting}

Now when we click on the link, we are re-addressed to the page of a specific product entity, that we can now modify to our liking.

\medskip

\textbf{A few words on JS Interop.}

In essence, JS Interop allows bidirectional communication between C\# and JavaScript, in other words, it allows C\# code to invoke JS functions, and JS functions to call C\# code.

Let's integrate JS Interop to our code. First, add folder named 'js' to wwwroot folder and add app.js file there. Write the following to it:

\begin{lstlisting}
function sayHello(name) {
    alert(`Hello ${name}`);
}
\end{lstlisting}

And add the following line to the App.razor file:

\begin{lstlisting}
<script src="js/app.js"></script>
\end{lstlisting}

Modify Prodcuts.razor file:

\begin{lstlisting}
@page "/products"
@using demo_sln.Models
@using demo_sln.Services
@rendermode InteractiveServer
@inject IProductsService _productsService
@inject IJSRuntime Js

<h3>Products</h3>

@if (_products == null)
{
    <h3>Please wait a bit</h3>
}
else
{
    <div>
        @foreach (var product in _products)
        {
            <ProductItem Product="product"></ProductItem>
        }
    </div>
}
<div>
    <input @bind="_name" />
    <input @bind="_description" />
    <button @onclick="SubmitData">Click me to add new product</button>
    <button @onclick="SayHello">Say Hello</button>
</div>



@code {
    private List<Product>? _products;

    private string? _name;
    private string? _description;

    protected override Task OnInitializedAsync()
    {
        _products = _productsService.GetProducts();
        return base.OnInitializedAsync();
    }

    private void SubmitData()
    {
        _productsService.AddProduct(
            new Product() { Name = _name, Description = _description }
        );
        _name = String.Empty;
        _description = String.Empty;
        _products = _productsService.GetProducts();
    }

    private async Task SayHello()
    {
        await Js.InvokeVoidAsync("sayHello", args: "products");
    }
}

\end{lstlisting}

Now you see browser alert any time you click button. You can integrate more Java Script functionality to your Razor components!

We can also call C\# code from JavaScript file using JSInterop. Take a look at the following static method in ProductsService.cs:

\begin{lstlisting}
[JSInvokable("LogMessage")]
public static Product LogMessage(int id)
{
    Console.WriteLine("Hello message " + id);
    return new Product() { Id = 3, Name = "test"};
}
\end{lstlisting}

We can call it from JavaScript file by specifying assembly (in our case, 'demo-sln'), method identifier (LogMessage) and parameter to be passed:

\begin{lstlisting}
async function callCSharpMethod() {
    const result = await DotNet.invokeMethodAsync("demo-sln", "LogMessage", 1);
    console.log(result);
}
\end{lstlisting}

Noiw by adding the following button:

\begin{lstlisting}
<button onclick="callCSharpMethod()">Call C# Method</button>
\end{lstlisting}

We can see the execution result being logged to console in both browser and DotNet console.


\medskip 

\textbf{Self-check questions:}

\begin{itemize}
    \item What is ASP.NET and what types of web applications can be built with it?
    
    \item What is Blazor and how does it differ from traditional JavaScript-based frontend frameworks?
    
    \item What is the difference between Blazor Server and Blazor WebAssembly hosting models?
    
    \item What is a Razor component and what are its main parts?
    
    \item What is the purpose of the \texttt{@page} directive in a Razor component?
    
    \item Why are interfaces used in C\# applications and what advantages do they provide?
    
    \item What is dependency injection and how is it implemented in ASP.NET applications?
    
    \item What is the purpose of the \texttt{@inject} directive in Blazor components?
    
    \item How does data binding work in Blazor and what is the role of the \texttt{@bind} directive?
    
    \item How can routing be implemented in Blazor applications using route parameters?
    
    \item What is JS Interop in Blazor and why is it important?
    
    \item How can C\# methods invoke JavaScript functions in Blazor applications?
    
    \item How can JavaScript code call C\# methods using JS Interop?
    
    \item What is the purpose of the \texttt{[JSInvokable]} attribute?
    
    \item Why is the Singleton lifetime useful in the demonstrated \texttt{ProductsService} example?
\end{itemize}

